\newpage
\section{Задание 1}

\subsection*{Постановка задачи}

Исходную IP-сеть необходимо разделить на 5 подсетей с заданным числом адресов в подсетях.

\noindent\textbf{Примечание:}
\begin{itemize}
    \item В задании указаны требования к числу адресов, а не к числу узлов.
    \item Подсетей должно быть 5, и они должны занимать всё адресное пространство исходной сети.
\end{itemize}

\begin{table}[H]
\centering
\begin{tabular}{|l|l|p{6.5cm}|}
\hline
\textbf{Исходная сеть} &
\textbf{Маска} &
\textbf{Требования к подсетям} \\ \hline
\texttt{168.60.20.0} &
\texttt{255.255.252.0} &
По крайней мере 2 подсети должны включать по 128 адресов
\\ \hline
\end{tabular}
\caption{Исходные данные (вариант 22)}
\label{tab:variant_data}
\end{table}

\subsection*{Ход выполнения}

Исходная сеть:

\[
\texttt{168.60.20.0/22}
\]

Маска сети:

\[
\texttt{255.255.252.0}
\]

Префикс \texttt{/22} означает, что для адресации сети используется 22 бита, а для адресации узлов остается:

\[
32 - 22 = 10 \text{ бит.}
\]

Поэтому исходная сеть содержит:

\[
2^{10} = 1024 \text{ адреса.}
\]


В маске \texttt{255.255.252.0} третий октет равен 252. Размер блока в третьем октете определяется как:

\[
256 - 252 = 4
\]

Следовательно, сети с такой маской начинаются в третьем октете с чисел, кратных 4. Число 20 кратно 4, поэтому \texttt{168.60.20.0} является адресом начала исходной сети.

Таким образом, исходная сеть охватывает диапазон:

\[
\texttt{168.60.20.0 -- 168.60.23.255}
\]

Следующая сеть с маской \texttt{/22} начнётся с адреса \texttt{168.60.24.0}, поэтому последний адрес исходной сети равен \texttt{168.60.23.255}.

Для получения 5 подсетей, полностью занимающих исходное адресное пространство, необходимо распределить 1024 адреса следующим образом:

\[
128 + 128 + 256 + 256 + 256 = 1024
\]

Подсеть из 128 адресов имеет префикс \texttt{/25}, поскольку:

\[
2^{32-25} = 128
\]

Подсеть из 256 адресов имеет префикс \texttt{/24}, поскольку:

\[
2^{32-24} = 256
\]

В результате получаем следующее распределение адресного пространства.

\begin{table}[H]
\centering
\begin{tabular}{|c|l|l|c|}
\hline
\textbf{Подсеть} &
\textbf{Адрес} &
\textbf{Маска} &
\textbf{Число адресов} \\ \hline

1 &
\texttt{168.60.20.0} &
\texttt{255.255.255.128} &
128 \\ \hline

2 &
\texttt{168.60.20.128} &
\texttt{255.255.255.128} &
128 \\ \hline

3 &
\texttt{168.60.21.0} &
\texttt{255.255.255.0} &
256 \\ \hline

4 &
\texttt{168.60.22.0} &
\texttt{255.255.255.0} &
256 \\ \hline

5 &
\texttt{168.60.23.0} &
\texttt{255.255.255.0} &
256 \\ \hline

\end{tabular}
\caption{Параметры подсетей}
\label{tab:subnets}
\end{table}

\noindent\textbf{Схема распределения адресов:}

\begin{itemize}
    \item \textbf{Подсеть 1}
    \begin{itemize}
        \item \texttt{168.60.20.0 -- 168.60.20.127}
        \item \texttt{/25} --- 128 адресов
    \end{itemize}

    \item \textbf{Подсеть 2}
    \begin{itemize}
        \item \texttt{168.60.20.128 -- 168.60.20.255}
        \item \texttt{/25} --- 128 адресов
    \end{itemize}

    \item \textbf{Подсеть 3}
    \begin{itemize}
        \item \texttt{168.60.21.0 -- 168.60.21.255}
        \item \texttt{/24} --- 256 адресов
    \end{itemize}

    \item \textbf{Подсеть 4}
    \begin{itemize}
        \item \texttt{168.60.22.0 -- 168.60.22.255}
        \item \texttt{/24} --- 256 адресов
    \end{itemize}

    \item \textbf{Подсеть 5}
    \begin{itemize}
        \item \texttt{168.60.23.0 -- 168.60.23.255}
        \item \texttt{/24} --- 256 адресов
    \end{itemize}
\end{itemize}

\noindent
Проверим, что все адресное пространство исходной сети распределено:

\[
128 + 128 + 256 + 256 + 256 = 1024
\]

Таким образом, получено ровно 5 подсетей, две из которых содержат по 128 адресов, а все 1024 адреса исходной сети \texttt{168.60.20.0/22} распределены между подсетями.